\documentclass[11pt]{article}
\usepackage{../homework}

\profname{Dr. David Singer}
\classcode{MATH 465 Differential Geometry}
\assignnum{Midterm, Take-Home Part}
\duedate{14 October 2026}

\begin{document}
\begin{problem}{1}
  Consider the planar curve \(\sigma(t) = (t\sin t + \cos t, -t\cos t +\sin t)\) for \( -2\pi \leq t \leq 8\pi \).
  \begin{enumerate}[label=(\alph*)]
    \item Reparametrize the curve by arc length.
          \begin{answer}
            The arc length is
            \begin{equation}
            \label{eq:3}
            \begin{aligned}
              s &= \int_{-2\pi}^{8\pi} | \sigma^{\prime}(t) | \dd{t} = \int_{-2\pi}^{8\pi} \left| (t\cos t + \sin t - \sin t, t\sin t - \cos t + \cos t) \right| \dd{t} \\
              &= \int_{-2\pi}^{8\pi} \left| (t\cos t, t\sin t) \right| \dd{t} = \int_{-2\pi}^{8\pi} |t| \dd{t} = \int\limits_{-2\pi}^0 -t \dd{t} + \int\limits_0^{8\pi} t \dd{t}.
            \end{aligned}
          \end{equation}
          Note that
          \begin{equation}
          \label{eq:5}
          \int\limits_{-2\pi}^0 -t \dd{t} = 2\pi^2,
          \end{equation}
          so that
          \begin{equation}
          \label{eq:4}
          \begin{aligned}
          s(t) &= \begin{cases}
            \int\limits_{-2\pi}^{t} -u \dd{u}, & -2\pi \leq t \leq 0 \\
            2\pi^2 + \int\limits_0^{t} u \dd{u}, & 0 \leq t \leq 8\pi
          \end{cases} \\
            &= \begin{cases}
              2\pi^2 - \frac{t^2}{2}, & -2\pi \leq t \leq 0 \\
              2\pi^2 + \frac{t^2}{2}, & 0 \leq t \leq 8\pi
            \end{cases}.
          \end{aligned}
        \end{equation}

        To find the inverse, we invert each piece separately, so that the first piece gives
        \begin{equation}
        \label{eq:6}
        \begin{aligned}
          s &= 2\pi^2 - t^2 / 2 \\
          \Rightarrow & t^2 = 4\pi^2 - 2s \\
          \Rightarrow & t = -\sqrt{4\pi^2 - 2s},
        \end{aligned}
      \end{equation}
      where we take the negative root because of the bounds, and the second gives
      \begin{equation}
      \label{eq:7}
      \begin{aligned}
        s &= 2\pi^2 + t^2 /2 \\
        \Rightarrow & t^2 = 2s - 4\pi^2 \\
        \Rightarrow & t = \sqrt{2s - 4\pi^2},
      \end{aligned}
    \end{equation}
    or
    \begin{equation}
    \label{eq:p1-tofs}
    t(s) = \begin{cases}
      -\sqrt{4\pi^2 - 2s}, & 0 \leq s \leq 2\pi^2 \\
      \sqrt{2s - 4\pi^2}, & 2\pi^2 \leq s \leq 34\pi^2
    \end{cases}.
  \end{equation}

  Define \( \alpha(s) = -\sqrt{4\pi^2 - 2s} \), \(\beta(s) = \sqrt{2s - 4\pi^2}\).
  Then, substituting \cref{eq:p1-tofs} into the equation for \(\sigma(t)\) gives
  \begin{equation}
  \label{eq:9}
  \sigma(t(s)) = \begin{cases}
    (\alpha(s)\sin\alpha(s) + \cos\alpha(s), -\alpha(s)\cos\alpha(s) + \sin\alpha(s)), & 0 \leq s \leq 2\pi^2 \\
    (\beta(s)\sin\beta(s) + \cos\beta(s), -\beta(s)\cos\beta(s) + \sin\beta(s)), & 2\pi^2 \leq s \leq 34\pi^2
  \end{cases}.
  \end{equation}
  Note that this is not differentiable everywhere, due to the misaligned left and right derivatives at \( s = 2\pi^2 \).
          \end{answer}
    \item Compute its Frenet frame \( (T(t), N(t)) \) and its (signed) curvature \( k(t) \) as functions of \( t \).
          \begin{answer}
            \( T(t) \) is, by definition, \(\sigma^{\prime}(t)/|\sigma^{\prime}(t)| = 1/|t|(t\cos t, t \sin t) \).
            Then \( N(t) = T^{\prime}/|T^{\prime}| = \frac{t}{|t|}(-\sin t, \cos t) \).

            Finally, the curvature is
            \begin{equation}
            \label{eq:17}
            k(t) = \frac{| T^{\prime}(t) |}{| \sigma^{\prime}(t) |} = \frac{1}{|t|} \left|(-\sin t, \cos t)\right| = \frac{1}{|t|}.
          \end{equation}

          All of this is for \( t \neq 0 \), since the frame and curvature are not defined at \( t = 0 \).
          \end{answer}
  \end{enumerate}
\end{problem}

\begin{problem}{2}\hfill
  \begin{enumerate}[label=(\alph*)]
    \item What does it mean to say that \( (T,U,V) \) is an inertial frame along a curve \(\sigma(s)\)?
          \begin{answer}
            An inertial frame is an orthonormal adapted frame \( (T, U, T \times U) \) where \( T \) is the unit tangent vector, \( \langle T, U \rangle = 0 \), and \( U \) is relatively parallel to \( T \), that is, \( U^{\prime}(s) = \gamma(s) T(s) \) for some scalar function \(\gamma(s)\).
          \end{answer}
    \item Suppose \(\sigma(s)\) is a unit speed curve lying on the unit sphere with \(\sigma^{\prime}(s) = T\).
          Let \( U(s) \) be the vector tangent to the sphere and orthogonal to \( T \) such that \( T \times U = V \) is the outward normal to the sphere.
          Verify that \( (T,U,V) \) is inertial.
          \begin{proof}
            It suffices to check that \( U \) is relatively parallel to \( T \).
            Since the frame is orthonormal by definition, we have that \(\langle U, U \rangle = 1\) and \( \langle U, V \rangle = 0 \).
            The first identity means that
            \begin{equation}
            \label{eq:8}
            \frac{\mathrm{d} }{\mathrm{d} s}\langle U, U \rangle = 2 \langle U^{\prime}, U \rangle = 0.
          \end{equation}

          Moreover, since \( \sigma \) lies on the surface of the sphere, we have \(\langle \sigma, \sigma \rangle = 1\), so that differentiation yields \( \langle \sigma, T \rangle = 0 \).
          Moreover, \( U \) is tangent to the sphere, so \(\langle \sigma, U \rangle = 0\).
          These two facts together with \( T \times U = V \) being the outward normal show that \(\sigma = V\), and \(\sigma^{\prime} = T = V^{\prime}\).
          We then differentiate the second identity from earlier, giving
          \begin{equation}
          \label{eq:11}
          \begin{aligned}
          \frac{\mathrm{d} }{\mathrm{d} s}\langle U, V \rangle &= \langle U^{\prime}, V \rangle + \langle U, V^{\prime} \rangle = \langle U^{\prime}, V \rangle + \langle U, T \rangle = \langle U^{\prime}, V \rangle = 0.
          \end{aligned}
        \end{equation}

        Using those two facts, we write \( U^{\prime} \) in the \( (T,U,V) \) basis as
        \begin{equation}
        \label{eq:12}
        U^{\prime} = a T + b U + c V,
      \end{equation}
      but since \( U^{\prime} \) is orthogonal to both \( U \) and \( V \), we must have that
      \begin{equation}
      \label{eq:13}
      U^{\prime} = a T,
    \end{equation}
    showing that \( U \) is relatively parallel to \( T \), and our claim follows.
          \end{proof}
    \item If \( T^{\prime}(s) = k_1U + k_2V \), express \( k_1 \) and \( k_2 \) in terms of the space curvature \( k \) of \(\sigma\).
          \begin{answer}
            Recall that, since \( \sigma \) is arc length parametrized, by definition \( | T^{\prime} | =  k\).
            Then, we have
            \begin{equation}
              \label{eq:curv-sqs}
              k^2 = k_1^2 + k_2^2.
            \end{equation}

            Since \( \langle T, V \rangle = 0 \),
            \begin{equation}
            \label{eq:14}
            \begin{aligned}
              \frac{\mathrm{d} }{\mathrm{d} s} \langle T, V \rangle &= \langle T^{\prime}, V \rangle + \langle T, V^{\prime} \rangle = \langle T^{\prime}, V \rangle + \langle T, T \rangle = k_2 + 1 = 0,
            \end{aligned}
          \end{equation}
          showing that \( k_2 = -1 \).
          It is immediate from \cref{eq:curv-sqs} that \( k_1 = \pm \sqrt{k^2 - 1} \).
          \end{answer}
  \end{enumerate}
\end{problem}

\begin{problem}{3}
  The surface \( M \) given by \( z = xy^2 \) is parametrized by
  \begin{equation}
  \label{eq:1}
  X(u,v) = (u,v,uv^2).
\end{equation}
\begin{enumerate}[label=(\alph*)]
  \item Compute the coefficients of the first fundamental form in \( (u,v) \) coordinates and the unit normal vector field \(\nu(u,v)\).
  \item Compute the coefficients of the second fundamental form of \( M \).
  \item Compute the Gaussian curvature of \( M \) at a general point \( X(u,v) \) and specifically at \( P= X(1,1) \) and \( Q=X(0,0) \).
\end{enumerate}
\end{problem}

\begin{problem}{4}
  Let \( M \) be a surface of revolution given by
  \begin{equation}
  \label{eq:2}
  X(z,\theta) = (f(z)\cos\theta, f(z)\sin\theta, z).
\end{equation}

\begin{enumerate}[label=(\alph*)]
  \item Show that any meridian is a line of curvature.
        (A meridian is a curve along which \(\theta\) is constant.)
  \item Why does it follow that any parallel is a line of curvature?
        (A parallel is a curve along which \( z \) is constant.)
\end{enumerate}
\end{problem}

\end{document}
